\section*{Devoir maison}

\begin{Exercise}[title={Somme des carrés}]

\begin{lstlisting}[style=pythonstyle, language=python]
x = int(input("Saisissez un nombre entier :"))
acc = 0
for i in range(x):
    acc += i**2
print(f"La somme des carres de 0 à {x} est {acc}")
\end{lstlisting}

\Question Rappelez l'utilité de la fonction input.
\Question Rappelez l'utilité de la fonction print.
\Question Rappelez l'utilité de la fonction range.
\Question Complétez le tableau d'avancement ci-dessous pour une saisie de x = 5.\\
\begin{center}
\large
\begin{tabular}{c|c|c|c|}
\textbf{Ligne} & \textbf{x} & \textbf{i} & \textbf{acc}  \\
\hline
 &  &  & \\
 &  &  & \\
 &  &  & \\
 &  &  & \\
 &  &  & \\
 &  &  & \\
 &  &  & \\
 &  &  & \\
\end{tabular}
\end{center}

\Question On souhaite maintenant modifier le programme pour calculer uniquement la somme des carrés pairs. Complétez le code ci-dessous pour obtenir ce comportement.
\begin{lstlisting}[style=pythonstyle, language=python]
x = int(input("Saisissez un nombre entier :"))
acc = 0
for i in range(x):
    ....................................................
        acc += i**2
print(f"La somme des carrés pairs de 0 à {x} est {acc}")
\end{lstlisting}

\Question On souhaite maintenant modifier le programme pour calculer la somme des carrés pairs  et la somme des carrés impairs entre un min et un max. Complétez le code ci-dessous pour obtenir ce comportement.
\begin{lstlisting}[style=pythonstyle, language=python]
xmin = int(input("Saisissez un nombre entier :"))
xmax = int(input("Saisissez un nombre entier :"))
acc_pair = 0
acc_impair = 0
for i in range(.............................):
    ....................................................
        acc_pair += i**2
    ....................................................
        acc_impair += i**2
print(f"La somme des carrés pairs de {xmin} à {xmax} est {acc_pair}, la somme des carrés impairs est {acc_impair}")
\end{lstlisting}
\end{Exercise}